% Clase 11 --- La corriente es un flujo (§7.5).
% Los dos argumentos con los que el docente lo justificó, uno por panel: al
% duplicar el área transversal pasa el doble de carga por unidad de tiempo
% (I proporcional a A), y con la superficie paralela al movimiento no la cruza
% nadie (I = 0). Eso es exactamente el comportamiento de un flujo: en lugar de
% contar líneas de campo se cuentan trayectorias de cargas.
\begin{center}
\begin{tikzpicture}[scale=0.95]

  % =============== (a) superficie normal al movimiento ===============
  \begin{scope}
    % trayectorias de las cargas
    \foreach \y in {-1.05,-0.7,...,1.1} {
      \draw[figvecaux, figred, line width=0.7pt,
            -{Latex[length=1.8mm,width=1.3mm]}] (-2.2,\y) -- (1.9,\y);
    }
    % superficie A (de canto)
    \draw[figblue, line width=1.8pt] (-0.5,-1.25) -- (-0.5,1.3);
    \node[figlbl, figblue, anchor=south] at (-0.5,1.35) {$A$};
    % normal
    \draw[figvec] (-0.5,-1.75) -- (0.35,-1.75);
    \node[figlbl, figblue, anchor=north] at (-0.05,-1.8) {$\hat n$};
    % superficie doble
    \draw[figgray, dashed, line width=1pt] (1.1,-2.1) -- (1.1,2.15);
    \node[figlblaux, anchor=south, align=center] at (1.1,2.2)
      {área doble\\[-0.2em] $\Rightarrow$ corriente doble};

    \node[figlbl, anchor=north, align=center] at (-0.4,-2.6)
      {(a) normal al movimiento\\[-0.2em] \footnotesize $I = J\,A$, y el
       cociente $I/A$ no depende del tamaño};
  \end{scope}

  % =============== (b) superficie paralela al movimiento ===============
  \begin{scope}[xshift=8.4cm]
    \foreach \y in {-1.05,-0.7,...,1.1} {
      \draw[figvecaux, figred, line width=0.7pt,
            -{Latex[length=1.8mm,width=1.3mm]}] (-2.0,\y) -- (2.0,\y);
    }
    % superficie contenida en el plano del movimiento
    \draw[figblue, line width=1.8pt] (-1.4,0.175) -- (1.4,0.175);
    \node[figlbl, figblue, anchor=south east] at (1.4,0.25) {$A$};
    \draw[figvec] (0,0.175) -- (0,1.0);
    \node[figlbl, figblue, anchor=west] at (0.08,0.7) {$\hat n$};
    \node[figlbl, figred, anchor=north] at (0,-1.5) {ninguna carga la cruza};

    \node[figlbl, anchor=north, align=center] at (0,-2.6)
      {(b) paralela al movimiento\\[-0.2em] \footnotesize
       $I = 0$: $\vec J \perp \hat n$};
  \end{scope}

\end{tikzpicture}\\[0.5em]
{\small\itshape Los dos comportamientos ---proporcional al área cuando la
superficie es transversal, nulo cuando es paralela--- son los mismos que ya
tenía el flujo del campo eléctrico, y por eso conviene describir la corriente
con la misma maquinaria: sólo cambia qué se cuenta, trayectorias de cargas en
vez de líneas de campo. Que $I$ sea proporcional a $A$ es justamente lo que
hace que el cociente $I/A$ \textbf{no dependa del tamaño} de la superficie
elegida, y por lo tanto que $\vec J$ sea una propiedad \emph{local} del punto y
no del área que se dibujó ahí. El sentido de $\vec J$ es el del movimiento
\emph{medio} de las cargas positivas ---o el opuesto al de las negativas, que
es el caso de un metal---, de modo que dos materiales con portadores de signo
contrario moviéndose en sentidos contrarios dan la \emph{misma} corriente.}
\end{center}
